\section{Scope, principal results, and conventions}
\label{intro:section}

The starting point is the ProveIt repository at revision
\texttt{03112a7d3b91} (full revision in \cite{ProveIt}), inspected on
12 October 2026. We read the polylogarithm manuscript and the text sources
in the incoming reports available at that revision; the most immediate
predecessors are the ordered-boundary, gap-six, and Gamma-normalization
reports \cite{ProveIt,OrderedIncoming,GapSix,GammaIncoming}.
This matters because several questions suggested by older manuscript
passages have already been answered in newer incoming material.
In particular, complete gap six, fixed-rate Stieltjes collisions, and
the paired Gamma deformation are prior results, not new contributions
of this article.

Our aim is to obtain exact evaluations and structural formulas. Estimates
are used where needed to prove identities or to turn eventual
nonsingularity into a theorem valid for all indices. They are not the
principal conclusions. Each of the three main directions gives a
finite, explicit way to evaluate an infinite family.

\subsection{Three advances and their exact scope}

\paragraph{Joint polygamma moments.}
For the positive $n$-variable Gamma density defined in
\eqref{gc:density}, introduce
$R_q=\sum_{j=1}^n\Re\psi^{(q)}(\sigma+ix_j/2)$.
Theorem~\ref{gc:main} evaluates all joint cumulants satisfying
$\sum_i(q_i+1)\le2n+2$. It gives means, covariances, and the particularly
simple higher-cumulant formula
\[
 \kappa(R_{q_1},\ldots,R_{q_r})
 =-2^{-r}\psi^{(r-1)}((2n+2)\sigma)
          \prod_i\delta_{q_i0},\qquad r\ge3.
\]
Thus the positive-order sector has exact Gaussian moments in a finite
weighted range, while the digamma sector retains a non-Gaussian
correction. This answers the joint-separation question of the latest
Gamma report in that range. The partition formula
\eqref{gc:moment} evaluates the corresponding ordinary integrals without
leaving parameter derivatives unevaluated. Proposition~\ref{gc:sharpness}
also proves a specific obstruction to extending the source's
single-trigamma matching degree universally.

\paragraph{Gap-eight reconstruction and sums.}
For $p=2b-2m$, $q=2b$, the harmonic-tail problem reduces to a rational
matrix $M_{b,m}$. Theorem~\ref{eg:effective} gives a fully computable
threshold $B(m)$ beyond which it is nonsingular, for every prescribed
$m\ge1$. Theorem~\ref{eg:eight-main} completes $m=4$ at every $b\ge5$,
settling the explicitly open gap-eight ray. Its finite step is an exact
computation, not a numerical extrapolation. Corollary~\ref{eg:eight-identities}
then gives rational polynomials $P_b,S_b$ such that
\[
 \sum_{n=0}^{\infty}\left[
 P_b(n+\tfrac12)\zeta(2b-8,n+1)\zeta(2b,n+1)
       -S_b(n+\tfrac12)\right]=-\frac12.
\]
The subtracted series is absolutely convergent. Matrix rank is taken on
the indicated formal coordinates; it is not a claim that the evaluated
zeta values are rationally or algebraically independent.

\paragraph{Different-power Stieltjes collisions.}
Theorem~\ref{mc:main} evaluates the change in every finite-part constant
when the two arguments are shifted by
$a=\lambda\epsilon^\alpha$, $b=\mu\epsilon^\beta$, with positive
parameters and unequal exponents. The equal-exponent case is inherited.
The answer consists of two explicit mixed contacts and at most one
singular--singular contact. The latter can survive precisely when
\[
 j(\alpha-\beta)=\beta(r+t+1),\qquad j\in\mathbb Z_{\ge1},
 \quad\alpha>\beta.
\]
The mixed contacts obey two further integer selection rules. This
classification explains constants absent from the fixed-rate theorem,
including the two renormalized limits $-1-\zeta(3)$ and $-1+\zeta(4)$
in \eqref{mc:q2limit}--\eqref{mc:q3limit}.
An especially short consequence is
\[
 \operatorname{CT}_{\epsilon}\int_0^1
 \gamma_1(x+\epsilon^{j+1})\gamma_0(x+\epsilon^j)\,dx
 =\frac32\gamma_2+2\gamma\zeta(2)-\zeta(3)+\psi'(j+1).
\]
Here the fixed finite part is inherited; the additional trigamma tail
is the newly identified collision term.

\subsection{Classical ingredients and the novelty boundary}

The multidimensional Gamma integral is a classical Gustafson integral;
we use the explicit symplectic normalization reproduced in
\cite[Eq.~(2)]{DerkachovManashov}, together with the repository's
determinant transformation. The new step is separation in independent
power-sum coordinates through their exact injectivity range.
The reconstruction proof builds on the repository's Bernoulli recurrence
and Gamma basis. Its new steps are a uniform effective threshold,
the complete gap-eight certificate, and the associated evaluated sums.

Finite-part integration and generalized Stieltjes-transform asymptotics
have an established literature \cite{TicaGalapon}; integrals of products
of Hurwitz zeta functions also occur in \cite{ShpotParis}. The monomial
collision theorem below uses elementary local subtraction, an incomplete
beta expansion, and cancellation of two meromorphic branches. Its new
content is the evaluated all-index contact law for the specified
two-scale problem. The Barnes companion in Section~\ref{bim:section}
is explicitly presented as a consequence of the classical first Barnes
lemma \cite[5.13.3]{DLMF}. We do not claim that parameter differentiation,
Bell polynomials, Barnes integration, or the resulting elementary
Euler-sum generators are new methods.

The comparison establishes advances over the inspected repository,
not an exhaustive historical-priority determination. The claims are
proved in the usual mathematical sense, with an explicitly isolated
exact computational component for the finite gap-eight interval.
The accompanying numerical calculations are independent diagnostics;
they are not proofs by integer-relation recognition. No proof-assistant
formalization is claimed.

\subsection{Notation and finite-part convention}

We use $\psi=\Gamma'/\Gamma$, $\psi^{(0)}=\psi$, and the standard
Hurwitz zeta function $\zeta(s,x)$. A superscript in parentheses on
$\gamma_m(x)$ denotes an argument derivative, whereas $m$ is its
Stieltjes index. The convention for generalized Stieltjes constants is
\[
 \zeta(1+u,x)=\frac1u+
       \sum_{m\ge0}\frac{(-1)^m}{m!}\gamma_m(x)u^m,
 \qquad \gamma_m=\gamma_m(1),\quad\gamma_0=\gamma.
\]
Thus $\gamma_0(x)=-\psi(x)$. Rising factorials are $(a)_k$;
falling factorials are explicitly underlined. We set
\[
 H_j^{(h)}=\sum_{k=1}^j k^{-h},\qquad H_j=H_j^{(1)},
 \qquad H_0^{(h)}=0.
\]
Bernoulli numbers have generating function
$t/(e^t-1)=\sum_{k\ge0}B_k t^k/k!$.
All logarithms of positive arguments are real; analytic logarithms
elsewhere are fixed locally at the stated base point.

The operation $\operatorname{CT}_{\epsilon}$ extracts the coefficient
of $\epsilon^0(\log\epsilon)^0$ from the specified power--logarithm
expansion. It is a coordinate-dependent finite part, not an unspecified
regularization and not necessarily a limit. We supply explicit ordinary
subtractions when stating a renormalized limit. The monomial exponents
are fixed while $\epsilon\downarrow0$; a jointly varying exponent ratio
near a resonance is a separate problem.
